\chapter{1985:Herbert A. Hauptman and Jerome Karle}

\label{chap:chemistry1985}

The Nobel Prize in Chemistry for the year 1985 was awarded jointly to Herbert A. Hauptman and Jerome Karle.

\textbf{Motivation:}

for their outstanding achievements in the development of direct methods for the determination of crystal structures

\section{Herbert A. Hauptman}

\subsection*{Early Life and Education}

Herbert A. Hauptman was born in 1917 in New York City, New York, USA.

\subsection*{Career and Major Research}

Hauptman studied at the City College of New York (B.S., 1937) and Columbia University (M.A., 1939), earning his doctorate in physics from the University of Maryland in 1955. He worked at the Naval Research Laboratory in Washington, D.C., and later led the Medical Foundation of Buffalo, which was renamed the Hauptman-Woodward Medical Research Institute in his honor.

\subsection*{Nobel Prize-winning Work}

The work for which Herbert A. Hauptman was awarded the Nobel Prize in Chemistry in 1985 was described by the Nobel Committee as: "for their outstanding achievements in the development of direct methods for the determination of crystal structures".

Together with Jerome Karle, he developed a mathematical framework for deriving the phases of X-ray reflections directly from the measured diffraction intensities, thereby solving the phase problem that had limited crystal structure analysis.

\subsection*{Significance and Impact}

Direct methods turned X-ray crystallography of small molecules into a routine and largely automated procedure, replacing the trial-and-error and heavy-atom approaches used before. The technique has been used to determine tens of thousands of crystal structures, including many important drugs and biological molecules.

\subsection*{Later Career and Honors}

Hauptman served as president of the Hauptman-Woodward Medical Research Institute in Buffalo for many years. He received the Nobel Prize in Chemistry in 1985 and died on 2011-10-23 in Buffalo, New York.

\subsection*{Legacy}

Direct methods remain the standard technique for small-molecule crystal structure determination, and the Hauptman-Woodward Medical Research Institute continues to bear his name.

\subsection*{Modern Developments}

Hauptman's direct methods for solving the phase problem in X-ray crystallography automated the determination of crystal structures, enabling tens of thousands of small-molecule structures to be solved each year. His mathematical approach underpins modern structural chemistry, drug development, and materials science and evolved into the phasing algorithms used in macromolecular crystallography.

\subsection*{Selected Publications}

"Solution of the Phase Problem I. The Centrosymmetric Crystal" (ACA Monograph, 1953)
"Crystal Structure Determination: The Role of the Cosine Seminvariants" (Plenum Press, 1972)

\clearpage

\section{Jerome Karle}

\subsection*{Early Life and Education}

Jerome Karle was born in 1918 in Brooklyn, New York, USA.

\subsection*{Career and Major Research}

Karle studied at the City College of New York (B.S., 1937) and Harvard University (A.M., 1938), earning his doctorate in physical chemistry from the University of Michigan in 1944. He worked at the Naval Research Laboratory in Washington, D.C., beginning in 1944, for his entire career.

\subsection*{Nobel Prize-winning Work}

The work for which Jerome Karle was awarded the Nobel Prize in Chemistry in 1985 was described by the Nobel Committee as: "for their outstanding achievements in the development of direct methods for the determination of crystal structures".

Working with Herbert A. Hauptman, he developed the statistical and mathematical methods that allowed the phases of X-ray reflections to be computed directly from measured intensities, solving the phase problem of X-ray crystallography.

\subsection*{Significance and Impact}

The direct methods developed by Hauptman and Karle made crystal structure determination faster and fully objective, and they were adopted by crystallographers worldwide. Karle also applied anomalous scattering of X-rays to determine the absolute configurations of molecules.

\subsection*{Later Career and Honors}

Karle remained at the Naval Research Laboratory until his retirement and continued research there. He received the Nobel Prize in Chemistry in 1985 and died on 2013-06-06 in Annandale, Virginia.

\subsection*{Legacy}

Karle's direct methods and his work on anomalous scattering remain foundational to modern X-ray crystallography, used in both chemistry and structural biology.

\subsection*{Modern Developments}

Karle's direct methods and his work on anomalous scattering enabled the automated solution of crystal structures and the determination of absolute configurations. These methods are used daily in structural chemistry and pharmaceutical development, while his diffraction techniques underlie modern crystallography of small molecules and of proteins.

\subsection*{Selected Publications}

"Solution of the Phase Problem I. The Centrosymmetric Crystal" (ACA Monograph, 1953)

\clearpage

\section*{Chapter Summary}

The 1985 Nobel Prize in Chemistry recognized the solution of the phase problem in X-ray crystallography. The direct methods developed by Hauptman and Karle turned structure determination from measured diffraction intensities into a routine, automated procedure used in laboratories worldwide.
